\documentclass[12pt, a4paper]{article}

%%------------------Preamble------------------%%

    %% 引入Package %%
\usepackage[margin=2cm]{geometry} % 上下左右距離邊緣3cm
\usepackage{fancyhdr}   % 頁首頁尾 fancy header
\usepackage{titling}    % 設定 title,author,date
\usepackage{mathtools,amsthm,amssymb} % 引入 AMS 數學環境
\usepackage[shortlabels,inline]{enumitem} %加強 enumerate, itemize 功能
\usepackage{setspace}
\usepackage{tcolorbox}
\usepackage{xcolor} % 用於表格上色
\usepackage{graphicx}      % 用於縮放表格寬度
\usepackage{array}
\usepackage{float} % for [H] float placement
\tcbuselibrary{breakable}
\usepackage{tabularx}
\usepackage{colortbl}

    %% 中文環境，僅限使用 XeLaTeX 編譯器 %%
\usepackage[CJKmath=true,AutoFakeBold=3,AutoFakeSlant=.2]{xeCJK}

\setCJKmainfont{Noto Serif CJK TC}
\setCJKsansfont{Noto Sans CJK TC}
\setCJKmonofont{Noto Sans Mono CJK TC}

\newCJKfontfamily\Kai{Noto Serif CJK TC}
\newCJKfontfamily\Hei{Noto Sans CJK TC}
\newCJKfontfamily\NewMing{Noto Serif CJK TC}

\XeTeXlinebreaklocale "zh"




\begin{document}

\title{T5 選擇題解答} % 設定標題
\author{李原廷} % 設定作者
\date{\today} % 設定完成日期，也可以使用 \date{\today}

\maketitle



\begin{enumerate}

    \item 下列關於超額負擔概念的敘述，何者正確？
    \begin{enumerate}[label=\Alph*.]
        \item 如被課稅商品的未受補償反應（uncompensated response）為零，即表示不會產生超額負擔
        \item 定額稅的課徵不會產生超額負擔，因此納稅人的福利水準不會下降
        \item 超額負擔為課稅引起的消費者剩餘的減少，加計生產者剩餘的減少之總和
        \item 在已存在扭曲性租稅的情況下，加課另一扭曲性租稅，所產生的整體超額負擔可能降低
    \end{enumerate}
    解答：
    \begin{enumerate}[label=\Alph*.]
        \item 超額負擔主要與受補償反應或替代效果有關，不是單純看未受補償需求反應。即使未受補償反應為零，也可能是所得效果與替代效果互相抵銷，仍可能有超額負擔。
        \item 定額稅不會扭曲相對價格，因此不產生超額負擔；但納稅人仍要繳稅，所得減少，所以福利水準仍會下降。
        \item 消費者剩餘與生產者剩餘的減少中，有一部分變成政府稅收，並非全是超額負擔。正確應為：$\text{超額負擔} = \Delta CS + \Delta PS - \text{稅收}$
        \item 在次佳理論下，若市場中已存在一項扭曲性租稅，再加課另一項租稅，有可能部分抵銷原本的相對價格扭曲，使整體資源配置反而改善，因此整體超額負擔可能降低。
    \end{enumerate}

    \item 在理論上欲藉選擇特定商品課貨物稅以達稅負累進效果，須被課稅商品具有下列何項條件？
    \begin{enumerate}[label=\Alph*.]
        \item 其所得彈性要大，而需求的價格彈性也要大
        \item 其所得彈性要大，但需求的價格彈性要小
        \item 其所得彈性要小，但需求的價格彈性要大
        \item 其所得彈性要小，而需求的價格彈性也要小
    \end{enumerate}
    解答：
    \begin{itemize}
        \item 所得彈性大：表示該商品多為高所得者消費，或所得愈高，消費比例愈高。例如奢侈品、高級消費品等。如此課稅後，稅負較容易落在高所得者身上，產生累進效果。
        \item 需求價格彈性小：表示即使價格因課稅而上升，消費者也不太會減少購買。如此可避免高所得者大量逃避或替代消費，使稅負較能真正落在該商品消費者身上。
    \end{itemize}


    \item 根據寇雷特與哈格法則（Corlett-Hague rule），現實生活中如果無法對休閒課稅，為使超額負擔最小，對休閒的互補品應如何課稅？
    \begin{enumerate}[label=\Alph*.]
        \item 輕課
        \item 重課
        \item 不課稅
        \item 與其他貨物課徵相同稅率
    \end{enumerate}
    解答：略。

    \item 假設 \(X\) 財貨的市場供給線是一條水平線，市場需求線是一條負斜率的直線。原先市場均衡價格是 5 元，均衡交易量是 8 單位。現若政府對該財貨的供給者課徵每單位稅額為 10 元的從量稅，假設稅後均衡交易量為 6 單位，試問以下何者正確？
    \begin{enumerate}[label=\Alph*.]
        \item 單位稅額 10 元完全由供給者負擔
        \item 政府稅收 90 元
        \item 超額負擔 10 元
        \item 生產者剩餘減損 70 元
    \end{enumerate}
    解答：
    \begin{enumerate}[label=\Alph*.]
    \item \textbf{錯誤。}
    因為供給線為水平線，表示供給完全彈性。課稅後，供給者仍取得原來的稅後價格：
    \[
    P_s = 5
    \]
    政府課徵每單位稅額：
    \[
    u = 10
    \]
    因此消費者支付價格為：
    \[
    P_c = P_s + u = 5 + 10 = 15
    \]
    所以單位稅額 \(10\) 元完全由消費者負擔，而非由供給者負擔。
    \[
    \text{供給者負擔} = 0,\qquad \text{消費者負擔} = 10
    \]
    \item \textbf{錯誤。}
    政府稅收等於每單位稅額乘以稅後交易量。已知稅後均衡交易量為：
    \[
    Q_1 = 6
    \]
    因此政府稅收為：
    \[
    TR = u \times Q_1
    \]
    \[
    TR = 10 \times 6 = 60
    \]
    所以政府稅收不是 \(90\) 元，而是：
    \[
    \boxed{60\text{ 元}}
    \]
    \item \textbf{正確。}
    超額負擔為課稅後交易量減少所造成的三角形損失。原均衡交易量為：
    \[
    Q_0 = 8
    \]
    稅後交易量為：
    \[
    Q_1 = 6
    \]
    交易量減少：
    \[
    Q_0 - Q_1 = 8 - 6 = 2
    \]
    因此超額負擔為：
    \[
    DWL = \frac{1}{2} \times u \times (Q_0 - Q_1)
    \]
    \[
    DWL = \frac{1}{2} \times 10 \times 2 = 10
    \]
    所以超額負擔為：
    \[
    \boxed{10\text{ 元}}
    \]
    \item \textbf{錯誤。}
    因為供給線為水平線，供給者在課稅前取得價格：
    \[
    P_0 = 5
    \]
    課稅後供給者取得的稅後價格仍為：
    \[
    P_s = 5
    \]
    因此供給者取得的單位價格沒有下降，生產者剩餘沒有減少：
    \[
    \Delta PS = 0
    \]
    所以生產者剩餘減損不是 \(70\) 元，而是0。
    \end{enumerate}

    \item 根據消費者的預算限制式，在所得固定的情況下，課徵定額稅（lump-sum tax）與下列何者不具有等價性？
    \begin{enumerate}[label=\Alph*.]
        \item 對所有財貨課徵單一稅率貨物稅
        \item 對所得課徵比例所得稅
        \item 對消費課徵一般消費稅
        \item 對休閒課徵從價稅
    \end{enumerate}
    解答：可以順便看第22題。
    \begin{enumerate}[label=\Alph*.]
    \item \textbf{具有等價性。}
    設原本預算限制式為：
    \[
    P_X X + P_Y Y = M
    \]
    若對所有財貨課徵單一稅率貨物稅 \(t\)，則預算限制式變為：
    \[
    (1+t)P_X X + (1+t)P_Y Y = M
    \]
    兩邊同除以 \(1+t\)：
    \[
    P_X X + P_Y Y = \frac{M}{1+t}
    \]
    這表示實質所得下降，但商品間相對價格不變，因此效果類似定額稅使預算線平行內移。
    所以此選項與定額稅具有等價性。
    \item \textbf{具有等價性。}
    若所得固定為 \(M\)，課徵比例所得稅，稅率為 \(t\)，則稅後所得為：
    \[
    (1-t)M
    \]
    預算限制式變為：
    \[
    P_X X + P_Y Y = (1-t)M
    \]
    因為所得是固定的，所以比例所得稅只是使可支配所得減少，不會改變商品間的相對價格。
    其效果等同於課徵某一金額的定額稅：
    \[
    T = tM
    \]
    因此：
    \[
    P_X X + P_Y Y = M - T
    \]
    所以此選項與定額稅具有等價性。
    \item \textbf{具有等價性。}
    一般消費稅若是對所有消費財課徵相同稅率 \(t\)，則預算限制式為：
    \[
    (1+t)P_X X + (1+t)P_Y Y = M
    \]
    化簡可得：
    \[
    P_X X + P_Y Y = \frac{M}{1+t}
    \]
    這與選項 A 相同，只是使消費者的實質購買力下降，並未改變財貨間的相對價格。
    因此，一般消費稅與定額稅具有等價性。
    \item \textbf{不具有等價性。}
    若考慮休閒 \(Z\)，令 \(w\) 為工資率，\(T\) 為總時間，則原本包含休閒的完整預算限制式可寫為：
    \[
    P_X X + P_Y Y + wZ = wT
    \]
    其中 \(wZ\) 表示休閒的機會成本。
    若對休閒課徵從價稅，稅率為 \(t\)，則休閒的價格變為：
    \[
    (1+t)w
    \]
    預算限制式變為：
    \[
    P_X X + P_Y Y + (1+t)wZ = wT
    \]
    此時休閒相對於消費財的價格改變，會影響消費者在消費與休閒之間的選擇。
    因此，對休閒課從價稅會改變相對價格，不只是使預算線平行移動，所以不等同於定額稅。
    所以此選項與定額稅不具有等價性。
    \end{enumerate}

    \item 若有 \(X\)、\(Y\) 兩種財貨，政府原已對 \(X\) 課稅，現又對 \(Y\) 課稅，則在 \(X\)、\(Y\) 兩財貨互為替代品或互補品兩種情況下，其所產生之總超額負擔將會如何？
    \begin{enumerate}[label=\Alph*.]
        \item 前者的總超額負擔一定增加，而後者不一定增加
        \item 前者的總超額負擔不一定增加，而後者一定增加
        \item 兩者的總超額負擔均會增加
        \item 兩者的總超額負擔不一定增加
    \end{enumerate}
    解答：
    \begin{enumerate}[label=\Alph*.]
    \item \textbf{錯誤。}
    \item \textbf{正確。}
    若 \(X\) 與 \(Y\) 為替代品，對 \(Y\) 課稅會使消費者轉向消費 \(X\)，因此：
    \[
    \frac{\partial X}{\partial P_Y} > 0
    \]
    由於 \(X\) 原先已被課稅，\(X\) 的消費量原本低於效率水準。現在因 \(Y\) 稅使 \(X\) 消費增加，可能減輕原本 \(X\) 稅造成的扭曲。
    因此，替代品情況下，新增 \(Y\) 稅所造成的總超額負擔不一定增加。
    若 \(X\) 與 \(Y\) 為互補品，對 \(Y\) 課稅會使 \(Y\) 消費下降，也會使 \(X\) 消費下降：
    \[
    \frac{\partial X}{\partial P_Y} < 0
    \]
    但 \(X\) 原本已被課稅，其消費量已低於效率水準。現在 \(X\) 消費又因 \(Y\) 稅而進一步下降，會使原有扭曲擴大。
    因此，互補品情況下，總超額負擔一定增加。
    \item \textbf{錯誤。}
    在替代品情況下，對 \(Y\) 課稅會使消費者增加對 \(X\) 的消費：
    \[
    \frac{\partial X}{\partial P_Y} > 0
    \]
    這可能減少原先 \(X\) 稅所造成的扭曲，因此總超額負擔不一定增加。
    \item \textbf{錯誤。}
    若 \(X\) 與 \(Y\) 為互補品，對 \(Y\) 課稅會使 \(Y\) 消費下降，並連帶使 \(X\) 消費下降：
    \[
    \frac{\partial X}{\partial P_Y} < 0
    \]
    由於 \(X\) 原本已被課稅，其消費量已低於效率水準，因此 \(X\) 消費進一步下降會使原有扭曲更嚴重。
    再加上對 \(Y\) 課稅本身也會造成新的扭曲，因此互補品情況下，總超額負擔會增加。
\end{enumerate}

    \item 租稅的超額負擔主要是隨下列何者而產生？
    \begin{enumerate}[label=\Alph*.]
        \item 所得效果
        \item 替代效果
        \item 財富效果
        \item 位移效果
    \end{enumerate}
    解答：租稅扭曲了經濟決策，導致相對價格的改變，發生替代效果，都會產生超額負擔。

    \item 社會福利成本（即為超額負擔）包括消費者福利損失與生產者福利損失。其中消費者福利損失與供需彈性之關係為：
    \begin{enumerate}[label=\Alph*.]
        \item 與需求彈性成反比，而與供給彈性成正比
        \item 與需求彈性成正比，而與供給彈性成反比
        \item 與需求彈性成正比，亦與供給彈性成正比
        \item 與需求彈性成反比，亦與供給彈性成反比
    \end{enumerate}
    解答：
    \begin{itemize}
        \item 超額負擔於需求彈性高低成正比。
        \item 超額負擔與供給彈性高低成正比。
        \item 可以用圖形分析。
    \end{itemize}

    \item 假定勞動者的行為決策只有 \(X\) 與 \(Y\) 兩財貨及休閒可供選擇，且無法對休閒課稅。其他情況不變，當工資率（休閒的價格）上升時，若 \(X\) 財貨的受補償需求隨之增加，而 \(Y\) 財貨的受補償需求隨之減少，則為使超額負擔極小化，對兩財貨之貨物稅應如何課徵？
    \begin{enumerate}[label=\Alph*.]
        \item 對 \(Y\) 課以較 \(X\) 為高的稅率
        \item 對 \(X\) 課以較 \(Y\) 為高的稅率
        \item 對兩者課以相同的稅率
        \item 對兩者課以相同的稅額
    \end{enumerate}
    解答：交叉彈性（Cross-Price Elacticity of Demand）

    當工資率上升時，也就是休閒的價格上升時：
    \[
    \frac{\partial h_X}{\partial w} > 0
    \]
    表示 \(X\) 財貨的受補償需求增加，所以 \(X\) 與休閒為替代品。
    又：
    \[
    \frac{\partial h_Y}{\partial w} < 0
    \]
    表示 \(Y\) 財貨的受補償需求減少，所以 \(Y\) 與休閒為互補品。
    因此，為了使超額負擔極小化，應對與休閒互補的 \(Y\) 財貨課較高稅率。

    \item 若受補償的勞動供給對工資率的彈性為 0.3，稅前工資率為每小時 \$200，勞動所得的稅率為 0.1，稅前勞動供給量為 5,000 小時，則對勞動課稅所造成的超額負擔（excess burden）為：
    \begin{enumerate}[label=\Alph*.]
        \item \$500
        \item \$1,000
        \item \$1,500
        \item \$2,000
    \end{enumerate}
    解答：
    已知：
	受補償勞動供給彈性：
    \[
    \varepsilon_L^c = 0.3
    \]
    稅前工資率：
    \[
    w = 200
    \]
    勞動所得稅率：
    \[
    t = 0.1
    \]
    稅前勞動供給量：
    \[
    L = 5,000
    \]
    勞動所得稅造成的每小時稅額為：
    \[
    tw = 0.1 \times 200 = 20
    \]
    受補償勞動供給量的變動為：
    \[
    \Delta L = \varepsilon_L^c \times L \times t
    \]
    \[
    \Delta L = 0.3 \times 5,000 \times 0.1 = 150
    \]
    超額負擔為稅楔造成的三角形面積：
    \[
    EB = \frac{1}{2} \times tw \times \Delta L
    \]
    \[
    EB = \frac{1}{2} \times 20 \times 150
    \]
    \[
    EB = 1,500
    \]

    \item 若 \(w\) 代表工資率，\(t\) 代表稅率，則個人所得稅與個人消費稅之課徵，會影響休閒相對價格如何改變？
    \begin{enumerate}[label=\Alph*.]
        \item 兩者均影響休閒相對價格降為 \(w \times (1-t)\)
        \item 兩者均影響休閒相對價格降為 \(w/(1+t)\)
        \item 前者影響休閒相對價格降為 \(w \times (1-t)\)，而後者影響休閒相對價格降為 \(w/(1+t)\)
        \item 前者影響休閒相對價格降為 \(w/(1+t)\)，而後者影響休閒相對價格降為 \(w \times (1-t)\)
    \end{enumerate}
    解答：跟第22題一並講解。

    \item 若對勞動所得課徵累進稅產生所得效果（\(IE\)）與替代效果（\(SE\)），則其超額負擔（\(EB\)）的相關敘述，下列何者正確？
    \begin{enumerate}[label=\Alph*.]
        \item \(IE\) 與 \(SE\) 均屬 \(EB\)，且由邊際稅率決定 \(EB\) 的大小
        \item 僅 \(SE\) 屬 \(EB\)，且由邊際稅率決定 \(EB\) 的大小
        \item \(IE\) 與 \(SE\) 均屬 \(EB\)，且由平均稅率決定 \(EB\) 的大小
        \item 僅 \(IE\) 屬 \(EB\)，且由平均稅率決定 \(EB\) 的大小
    \end{enumerate}
    解答：
    \begin{itemize}
        \item 所得效果 \(IE\)：課稅使可支配所得下降，若休閒為正常財，所得下降會使勞動者減少休閒、增加工作。所得效果本身只是所得被政府拿走所造成的福利下降，類似定額稅效果，不屬於超額負擔。
        \item 替代效果 \(SE\)：所得稅降低稅後工資，使休閒相對變便宜，勞動者會以休閒替代工作。這會扭曲勞動與休閒的選擇，因此會產生超額負擔。
    \end{itemize}
    勞動者是否多工作一小時，取決於多工作一小時可得到的稅後邊際工資。
    若稅前工資為 \(w\)，邊際稅率為 \(t_m\)，則稅後邊際工資為：
    \[
    (1-t_m)w
    \]
    因此，影響勞動與休閒選擇的是邊際稅率。

    \item 下列何者非屬雷姆西法則（Ramsey rule）之應用？
    \begin{enumerate}[label=\Alph*.]
        \item 對於非主要工作者的所得課以較輕之租稅負擔
        \item 對於所得較高的家庭課以較高之租稅負擔
        \item 對於政府提供財貨中受補償需求彈性較大者索取較低的使用費
        \item 對於民生必需品課以較高的稅率
    \end{enumerate}
    解答：
    \begin{enumerate}[label=\Alph*.]
        \item 非主要工作者，例如家庭中的第二所得者，其勞動供給彈性通常較大。依雷姆西法則，彈性較大者應課較低稅率，以減少勞動供給扭曲。
        \item 這是基於量能課稅原則或所得重分配考量，屬於公平面或累進稅制的概念，不是雷姆西法則。雷姆西法則主要關心的是效率，也就是如何降低超額負擔。
        \item 需求彈性較大者，若收費太高，需求量會大幅減少，造成較大效率損失，因此應收取較低使用費。
        \item 民生必需品需求彈性通常較小，依反彈性法則，若只考慮效率，應課較高稅率。雖然這可能有公平問題，但就雷姆西法則而言是其應用。
    \end{enumerate}

    \item 下列何者非屬最適租稅問題的一般設定？
    \begin{enumerate}[label=\Alph*.]
        \item 為取得一定之稅收
        \item 所有租稅皆具扭曲性
        \item 稅收之支用所造成扭曲的極小
        \item 課稅為使社會福利的極大
    \end{enumerate}
    解答：
    最適租稅理論（optimal taxation）的一般設定是：
    \begin{itemize}
        \item 政府需要取得一定稅收。
        \item 因定額稅通常不可行或受限制，所以必須使用扭曲性租稅。
        \item 政府在課稅造成的效率損失與公平目標之間取捨。
        \item 目標是在既定稅收需求下，使社會福利極大，或使超額負擔極小。
    \end{itemize}
    可表示為：
    \[
    \max W
    \]
    subject to：
    \[
    R \geq \bar{R}
    \]
    其中 \(W\) 是社會福利，\(R\) 是政府稅收，\(\bar{R}\) 是政府所需稅收。
    \begin{enumerate}[label=\Alph*.]
        \item 正確。最適租稅通常是在政府需要籌措一定稅收的前提下，尋找最不扭曲、或最能兼顧公平與效率的稅制。
        \item 正確。若可使用無扭曲的定額稅，問題會變得簡單。因此最適租稅理論通常假設政府無法完全使用定額稅，而必須使用具有扭曲性的租稅工具。
        \item 錯誤。這是政府支出面或公共支出效率的問題。最適租稅主要討論的是「如何課稅」以降低課稅造成的扭曲，而不是討論稅收支用本身造成的扭曲。
        \item 正確。最適租稅的核心目標就是在政府稅收需求與租稅扭曲限制下，使社會福利極大。
    \end{enumerate}


    \item 假設只有 \(X\)、\(Y\) 和休閒三個財貨，政府只對 \(X\) 和 \(Y\) 課稅，無法對休閒課稅。下列敘述何者正確？
    \begin{enumerate}[label=\Alph*.]
        \item 依 Haig-Simons 的所得定義，一納稅人的所有所得來源應被課相同稅率，可使課稅產生的超額負擔最小
        \item 中性課稅必定符合 Ramsey 的課稅法則
        \item 為使課稅產生的超額負擔最小，稅率的設計應使 \(X\) 和 \(Y\) 減少相同的需求量
        \item 若 \(X\) 和 \(Y\) 既非替代亦非互補，且 \(X\) 和 \(Y\) 的受補償需求的價格彈性相同，則中性課稅（neutral taxation）能使課稅產生的超額負擔最小
    \end{enumerate}
    解答：
    \begin{enumerate}[label=\Alph*.]
        \item 錯誤。Haig-Simons 所得定義強調綜合所得概念，主張所得來源應完整納入稅基，但這主要是所得稅公平與稅基完整性的概念，不是使超額負擔最小的 Ramsey 法則。
        \item 錯誤。中性課稅不一定符合 Ramsey 法則。若 \(X\)、\(Y\) 的受補償需求彈性不同，則依 Ramsey 法則，應對彈性小者課較高稅率、彈性大者課較低稅率。此時相同稅率反而不是最適。
        \item 錯誤。Ramsey 法則不是要求 \(X\)、\(Y\) 減少「相同的需求量」。在特定簡化情況下，是要求受補償需求量的比例變動相同，而不是絕對數量變動相同。
        \item 若 \(X\) 和 \(Y\) 既非替代亦非互補，表示兩者間的受補償交叉價格效果為零。又若兩者受補償需求價格彈性相同，則依 Ramsey 反彈性法則：
        \[
        \varepsilon_X^c = \varepsilon_Y^c
        \]
        因此最適稅率為：$t_X = t_Y$。
        此時中性課稅可以使超額負擔最小。
    \end{enumerate}

    \item 為取得一定稅收，其他條件不變時，對單一財貨課稅，效率損失高於對數個財貨同時課稅，原因在於：
    \begin{enumerate}[label=\Alph*.]
        \item 無謂損失隨稅率增加之平方增加
        \item 對數個財貨同時課稅時，不同替代效果相互抵消作用
        \item 對數個財貨同時課稅時，不同彈性相互抵消作用
        \item 無謂損失隨被課稅貨物數量變化之平方增加
    \end{enumerate}
    解答：略。

    \item 在只有兩種被課稅財貨的情形下，下列有關雷姆西法則（Ramsey rule）的敘述何者正確？
    \begin{enumerate}[label=\Alph*.]
        \item 課稅所造成兩財貨之受補償需求量變動相同
        \item 最後一元稅收所造成兩財貨之邊際超額負擔相同
        \item 課稅所造成兩財貨之受補償需求量變動與財貨價格比例相同
        \item 稅率須使兩財貨之邊際替代率相同
    \end{enumerate}
    解答：
    \begin{enumerate}[label=\Alph*.]
        \item 雷姆西法則不是要求兩財貨的受補償需求量「絕對變動量」相同，而是要求在最適租稅下，邊際超額負擔相等。若在簡化情況下，也常表達為各財貨受補償需求量的「比例變動」相同，而不是絕對數量變動相同。
        \item 雷姆西法則的邊際條件。
        \item 不是要求受補償需求量變動與價格比例相同。
        \item 邊際替代率是消費者偏好的概念，不是雷姆西最適商品稅率的判準。
    \end{enumerate}

    \item 下列超額負擔（excess burden）的討論，何者錯誤？
    \begin{enumerate}[label=\Alph*.]
        \item 又被稱之為福利成本（welfare cost）
        \item 也被稱之為無謂損失（deadweight loss）
        \item 就課稅而言，所指的是稅收超過以金錢衡量之福利損失的額度
        \item 就課稅而言，所指的是相等變量（equivalent variation）絕對值大於稅收的部分
    \end{enumerate}
    解答：
    超額負擔（excess burden）是指課稅造成的福利損失中，超過政府稅收的部分。
    也就是：
    \[
    EB = \text{福利損失} - \text{稅收}
    \]
    若用相等變量（equivalent variation, EV）衡量福利損失，則可寫為：
    \[
    EB = |EV| - T
    \]
    其中：\(|EV|\)：納稅人因課稅所受的金錢化福利損失、\(T\)：政府稅收、\(EB\)：超額負擔
    \begin{enumerate}[label=\Alph*.]
        \item 正確。
        \item 正確。
        \item 方向相反。超額負擔不是「稅收超過福利損失」，而是「福利損失超過稅收」的部分。
        \item 正確。
    \end{enumerate}

    \item 考慮休閒、勞動可變動下，對勞動所得課徵比例所得稅（\(IT\)），或是對所有消費（不含休閒）課徵比例消費稅（\(CT\)），若其他條件相同，下列敘述何者正確？
    \begin{enumerate}[label=\Alph*.]
        \item \(IT\) 具勞動市場效率，\(CT\) 不具勞動市場效率
        \item \(IT\) 與 \(CT\) 均具勞動市場效率
        \item \(IT\) 不具勞動市場效率，\(CT\) 具勞動市場效率
        \item \(IT\) 與 \(CT\) 均不具勞動市場效率
    \end{enumerate}
    解答：略，第22題。

    \item 假設某稅制固定稅率為 20\%。逃漏稅被稽查成功除應補徵稅負外，另應納罰款，罰款公式為 \(F = 2y^2 + y\)，\(y\) 為逃漏之應稅所得金額，而逃漏稅被成功查獲之機率為 10\%。假設所得為 1 元，則最適逃稅金額為多少元？
    \begin{enumerate}[label=\Alph*.]
        \item 0.1 元
        \item 0.2 元
        \item 0.4 元
        \item 0.8 元
    \end{enumerate}
    解答：
    稅率為：
    \[
    t = 20\% = 0.2
    \]
    查獲機率為：
    \[
    p = 10\% = 0.1
    \]
    逃漏所得金額為 \(y\)，若未被查獲，可少繳稅額為：
    \[
    ty = 0.2y
    \]
    但只有未被查獲時才真正保有逃稅利益，未被查獲機率為：
    \[
    1-p = 0.9
    \]
    因此逃稅的預期利益為：
    \[
    (1-p)ty = 0.9 \times 0.2y = 0.18y
    \]
    若被查獲，除補稅外，另須繳罰款：
    \[
    F = 2y^2 + y
    \]
    其預期罰款成本為：
    \[
    pF = 0.1(2y^2+y)
    \]
    \[
    =0.2y^2+0.1y
    \]
    所以逃稅的預期淨利益為：
    \[
    E(G)=0.18y-(0.2y^2+0.1y)
    \]
    \[
    E(G)=0.08y-0.2y^2
    \]
    要求最適逃稅金額，對 \(y\) 微分：
    \[
    \frac{dE(G)}{dy}=0.08-0.4y
    \]
    令其等於 0：
    \[
    0.08-0.4y=0
    \]
    \[
    0.4y=0.08
    \]
    \[
    y=0.2
    \]
    因此最適逃稅金額為：
    \[
    \boxed{0.2\text{ 元}}
    \]

    \item 下列何者符合均等變量（equivalent variation）之定義？
    \begin{enumerate}[label=\Alph*.]
        \item 納稅人最多願意支付多少代價換取政府不要課稅的金額
        \item 對資本要素課稅後，因要素替代效果而使勞動要素價格上漲的金額
        \item 納稅人因繳稅所犧牲的邊際效用相等之金額
        \item 稅法為達成特定目的給予租稅減免，稅收損失之金額
    \end{enumerate}
    解答：
    均等變量（equivalent variation, EV）是衡量價格變動或課稅所造成福利損失的一種方法。
    若政府課稅使納稅人的效用從原來的 \(U^0\) 下降到 \(U^1\)，則均等變量是指：
    在課稅尚未發生前，納稅人最多願意支付多少金額，以避免政府課稅。
    也就是說，納稅人願意付出一筆錢，使自己在「不課稅但支付該金額」的情況下，效用剛好等於「政府課稅後」的效用。
    可表示為：
    \[
    EV = e(P^0,U^0) - e(P^0,U^1)
    \]
    其中：\(P^0\)：課稅前價格、\(U^0\)：課稅前效用、\(U^1\)：課稅後效用、\(e(P,U)\)：支出函數
    因此，選項 A 符合均等變量的定義。
    \begin{enumerate}[label=\Alph*.]
        \item 均等變量的概念。
        \item 這是要素稅轉嫁或一般均衡效果，不是均等變量。
        \item 這較接近租稅公平中的「等犧牲原則」，不是均等變量。
        \item 這是「租稅支出」或「稅式支出」的概念，不是均等變量。
    \end{enumerate}

    \item 在考慮休閒選擇可變動的條件下，課徵定額稅與下列何者具有等價性？
    \begin{enumerate}[label=\Alph*.]
        \item 對所得課徵比例所得稅
        \item 對休閒課徵從價稅
        \item 對所有財貨（不含休閒）課徵單一稅率的貨物稅
        \item 對休閒與所有財貨課徵單一稅率的貨物稅
    \end{enumerate}
    解答：可再講解第11題。

    若考慮休閒：
	\(Z\)：休閒、\(T\)：總時間、\(L = T - Z\)：勞動時間、\(w\)：工資率、\(P_X, P_Y\)：商品價格
    原本預算限制式為：
    \[
    P_X X + P_Y Y = wL
    \]
    因為：
    \[
    L = T - Z
    \]
    所以：
    \[
    P_X X + P_Y Y = w(T-Z)
    \]
    整理得：
    \[
    P_X X + P_Y Y + wZ = wT
    \]
    其中 \(w\) 是休閒的機會成本。若課定額稅 \(T_0\)，則：
    \[
    P_X X + P_Y Y + wZ = wT - T_0
    \]
    定額稅只會使全所得下降，沒有改變相對價格。

    （A）：若課比例所得稅，稅率為 \(t\)，則稅後工資變成：
    \[
    (1-t)w
    \]
    預算限制式為：
    \[
    P_X X + P_Y Y = (1-t)w(T-Z)
    \]
    整理得：
    \[
    P_X X + P_Y Y + (1-t)wZ = (1-t)wT
    \]
    此時休閒的價格從 \(w\) 變成：
    \[
    (1-t)w
    \]
    所以比例所得稅會降低休閒的機會成本，使休閒相對變便宜，進而影響勞動與休閒選擇。
    因此在休閒可變動時，比例所得稅會扭曲選擇，不等同於定額稅。

    （B）：對休閒課徵從價稅，稅率為 \(t\)，則休閒價格為：
    \[
    (1+t)w
    \]
    預算限制式為：
    \[
    P_X X + P_Y Y + (1+t)wZ = wT
    \]
    此時休閒的價格從 \(w\) 變成：
    \[
    (1+t)w
    \]
    所以從價稅會增加休閒的機會成本，使休閒相對變昂貴，進而影響勞動與休閒選擇。

    （C）：對所有財貨（不含休閒）課徵單一稅率的貨物稅，稅率為 \(t\)，則財貨價格為：
    \[
    (1+t)P_X, (1+t)P_Y
    \]
    預算限制式為：
    \[
    (1+t)P_X X + (1+t)P_Y Y + wZ = wT
    \]
    也可以寫為
    \[
    P_X X + P_Y Y + \frac{w}{(1+t)}Z = wT
    \]
    此時財貨的價格從 1 變成：
    \[
    (1+t)P_X, (1+t)P_Y
    \]
    也改變了財貨跟休閒的相對價格，影響財貨與休閒選擇。

    \item 哈伯格（A. C. Harberger）以一般均衡分析法，探討對公司部門的資本課稅的轉嫁與歸宿，下列敘述何者錯誤？
    \begin{enumerate}[label=\Alph*.]
        \item 可以從產出效果與要素替代效果分析其影響
        \item 若公司部門是資本密集產業，則透過產量效果與要素替代效果會將租稅部分轉嫁給非公司部門
        \item 可能會使勞動相對價格下降
        \item 若公司部門是勞動密集產業，且生產要素可以相互替代，稅負完全由公司部門的勞動來承擔
    \end{enumerate}
    解答：
    （D）：在哈伯格模型中，勞動與資本都可以在部門間自由移動，因此公司部門勞動不可能單獨完全承擔稅負。即使公司部門是勞動密集產業，稅負也可能分散到：公司部門勞動、非公司部門勞動、公司部門資本等。而且因為存在要素替代效果，公司部門會以勞動替代資本，這也會影響資本報酬率，因此不能說稅負「完全」由公司部門勞動承擔。

    \item 拉弗曲線（Laffer curve）意涵當所得稅率降低，且休閒為正常財，則稅收增加的條件是：
    \begin{enumerate}[label=\Alph*.]
        \item 替代效果小於所得效果，使工作時數隨著淨工資減少而增加
        \item 替代效果大於所得效果，使工作時數隨著淨工資減少而增加
        \item 替代效果小於所得效果，使工作時數隨著淨工資提高而增加
        \item 替代效果大於所得效果，使工作時數隨著淨工資提高而增加
    \end{enumerate}
    解答：
    所得稅率降低時，勞動者的稅後工資，也就是淨工資會上升：
    \[
    w_n = (1-t)w
    \]
    當所得稅率 \(t\) 降低時：
    \[
    w_n \uparrow
    \]
    若休閒為正常財，淨工資上升會產生兩種效果：
    \begin{itemize}
        \item 替代效果：淨工資提高，使休閒的機會成本上升，因此勞動者會減少休閒、增加工作時數。
        \item 所得效果：淨工資提高，使勞動者實質所得增加。因休閒是正常財，所以勞動者會增加休閒、減少工作時數。
    \end{itemize}
    若要所得稅率降低後，政府稅收反而增加，勞動供給必須增加得夠多，使稅基擴大足以抵銷稅率下降的效果。
    因此方向是：
    \[
    \text{替代效果} > \text{所得效果}
    \]
    使得：
    \[
    w_n \uparrow \quad \Rightarrow \quad L \uparrow
    \]
    也就是淨工資提高時，工作時數增加。

    \item 下列何者非屬平衡預算租稅歸宿分析之研究情況？
    \begin{enumerate}[label=\Alph*.]
        \item 在支出不變的情形下，提高所得稅同時降低消費稅
        \item 提高所得稅並以其稅收用之於愛滋病之防治
        \item 提高消費稅並以其稅收用之於提供弱勢家庭國民小學學童免費營養午餐
        \item 提高消費稅將所增稅收以定額方式（as a lump sum）退還納稅義務人使用
    \end{enumerate}
    解答：平衡預算租稅歸宿分析是指：政府增加某項租稅，同時將增加的稅收用於某項公共支出或移轉支出，使政府預算維持平衡。
    \begin{enumerate}[label=\Alph*.]
        \item 這是在政府支出不變下，用一種稅取代另一種稅，屬於差別租稅歸宿分析，不是平衡預算歸宿分析。
        \item 因為稅收增加，同時政府支出也增加。
        \item 因為增加的稅收用於特定公共支出或移轉支出。
        \item 因為政府增加稅收後，又以定額移轉方式退還，仍是在分析「增稅與支出或移轉」的合併效果。
    \end{enumerate}



\end{enumerate}




\end{document}
